\section{Conclusions : résolution affirmative de l’hypothèse du continu}
\label{sec:conclusiones-ch}
\begingroup
\fontsize{10.5}{13.4}\selectfont
\setlength{\parskip}{3pt}
\setlength{\abovedisplayskip}{5pt}
\setlength{\belowdisplayskip}{5pt}

Le résultat principal est l’absence de cardinaux intermédiaires entre les
naturels et la droite dans la clôture généalogique HMT. Pour chaque histoire
réalisée, la construction de cette clôture et les deux injections cardinales
établissent
\begin{equation}
 \boxed{
 \mathfrak G_{\mathrm{HMT}}(h,m_\infty)
 \models 2^{\aleph_0}=\aleph_1 .}
 \label{eq:conclusion-decision-continuo}
\end{equation}
L’égalité n’est présupposée à aucune étape générative. La couverture
exhaustive des cylindres établit l’équipotence
\[
 \mathbb R_{\mathrm{HMT}}^{\mathfrak G}
 \approx\mathcal P^{\mathfrak G}(\omega).
\]
La condensation généalogique attribue à chaque réel interne un ordinal strictement inférieur à \(\omega_1^{\mathfrak G}\) et construit
\[
 j_{\mathrm{HMT}}:
 \mathbb R_{\mathrm{HMT}}^{\mathfrak G}\hookrightarrow
 \omega_1^{\mathfrak G}.
\]
L’encodage des bons ordres dénombrables et leur relèvement ternaire
construisent l’injection dans le sens contraire
\[
 k_{\mathrm{HMT}}:
 \omega_1^{\mathfrak G}\hookrightarrow
 \mathbb R_{\mathrm{HMT}}^{\mathfrak G}.
\]
Le théorème de Cantor--Bernstein conclut l'égalité cardinale de \eqref{eq:conclusion-decision-continuo} ; il ne la fournit pas en entrée.

Le quantificateur du résultat porte sur tous les sous-ensembles de la droite qui
appartiennent à la clôture généalogique complète. Pour chaque
\(A\in\mathcal P^{\mathfrak G}
(\mathbb R_{\mathrm{HMT}}^{\mathfrak G})\), la cofinalité démontrée impose
la dichotomie
\[
 |A|\leq\aleph_0
 \qquad\text{ou}\qquad
 |A|=|\mathbb R_{\mathrm{HMT}}^{\mathfrak G}|.
\]
Il n’existe pas, dans le continu intégral engendré par HMT, de cardinal strictement
intermédiaire entre \(\aleph_0\) et le cardinal de la droite. Par conséquent,
l’holographie modulaire triadique résout affirmativement l’hypothèse du
continu dans la clôture généalogique complète définie dans ce travail.

La pertinence structurelle de cette égalité réside dans son lien avec une construction explicite de la droite. La composition
\[
 \mathrm{APP}\longrightarrow\mathrm{TRIT}\longrightarrow\mathrm{TPK}
 \longrightarrow X_\infty^{\mathrm{enr}}
 \longrightarrow\mathcal C_{\mathrm{cont}}^{\mathrm{disc}}
\]
engendre un unique système compatible d'histoires dotées de valeur, de feuillet, d'orientation, de résidu, de quotient, de retenue, de frontière, de chemin, de mémoire et de survie. Sa projection archimédienne produit la droite \(\mathbb R_{\mathrm{HMT}}^{\mathfrak G}\) ; sa projection d'incidence conserve l'information de frontière que l'évaluation de la coordonnée archimédienne ne retient pas. La droite réelle est une sortie de la construction, tandis que l'état conserve les opérations et les relations dont provient cette sortie.

La finitude de chaque niveau ne restreint pas cette portée. La fibre
\(\mathbb F_3^6\), formée de \(3^6=729\) sous-cylindres de raffinement,
décrit la ramification locale ; le quantificateur coinductif agit à toute
profondeur prescrite. Les exécutions à mille, dix mille ou un million de chiffres
sont des témoins finis d’une loi paramétrée par la profondeur, et non des bornes de
la construction. Le retour nonadique retrouve la phase tandis que la retenue et la
mémoire avancent, de sorte que la fermeture observable ne réinitialise pas l’état
enrichi.

% BEGIN R27_FRAMING_CONCLUSION_ES
Le développement de l’émetteur précise trois niveaux de conservation. La
congruence de contrôle du troisième chiffre accompagne les deux signatures
récupérables et détermine un alphabet de \(42\) blocs atteignables.
Les six émissions témoins construisent un cycle d’incidences régionales
dont la provenance dans les germes est explicite. La comparaison de deux mots
ayant les mêmes fréquences et une corrélation temporelle différente identifie
l’information d’ordre que conserve l’archive. Leurs cardinaux, graphes
et fonctionnelles appartiennent à des domaines finis définis ; la preuve de
l’affirmation cardinale principale conserve séparément ses quantificateurs,
sa clôture et les deux injections construites.
% END R27_FRAMING_CONCLUSION_ES

% BEGIN R28_FRAMING_CONCLUSION_ES
Les cinq clôtures régionales atteignent le minimum huit dans le problème fini de couverture des quatre arêtes prescrites. Le graphe conserve les étiquettes des émissions et le quotient conserve leurs incidences entre classes de phase. L'inversion orientée rend impair le lecteur \(\nu_{120}\) et maintient pair \(\nu_{270}\) ; ces identités explicitent la réponse de la mémoire au changement d'orientation.
% END R28_FRAMING_CONCLUSION_ES

\clearpage
La même généalogie produit les résultats corrélés de la double réalisation. Parmi eux figurent l'holonomie nonadique \(\Gamma_9\), sa lecture ultérieure de mémoire réduite \(\mathcal K_{\mathrm{ph}}\), la constante holographique de couplage arithmético-géométrique \(K\), la quadrirelation \((\pi,\varphi,e,\alpha)\) et le corridor Paley--Witt--Golay--Leech--\(V^\natural\)--Monstre. Ce sont des codomaines et des lecteurs distincts d'un même état ; aucun n'est introduit à partir de valeurs conventionnelles pour forcer la comparaison cardinale.

En particulier, \(K\) est la représentation entière canonique du registre
dodécaphasique orienté. Sa production suit la chaîne
\[
 X_{\rm term}\longrightarrow\mathcal R_{12}(X_{\rm term})
 \longrightarrow(b^{90},b^{120},Q_{\rm TPK})
 \longrightarrow U_{\rm sgn},
 \qquad
 U_{\rm sgn}\xrightarrow{\;\mathcal H_{12}/4\;}K,
 \qquad
 K\xrightarrow{\;\mathcal H_{12}\;}U_{\rm sgn}.
\]
L’équivalence de Hadamard exprime une mise en coordonnées réversible du
registre engendré ; elle ne transforme ni \(U_{\rm sgn}\), ni \(K\), ni sa lecture
rationnelle en données externes. La lecture positionnelle de \(K\) intervient
ultérieurement dans la clôture de \(\alpha\) ; ses lectures incidencielle et
géométrique conservent le support, la direction et la réduction dimensionnelle. Cette fonction
ne doit pas être confondue avec \(\mathcal K_{\mathrm{ph}}\), qui appartient à un autre
type et s’obtient par réduction de mémoire après l’holonomie.

L’identification
\[
 \mathfrak G_{\mathrm{HMT}}(h,m_\infty)
 =L[u_{h,m}]=L[h,m_\infty]
\]
est la reconnaissance finale d’une hiérarchie déjà construite. Elle explicite sa
représentation par constructibilité relative, mais n’engendre pas la droite, ne
sélectionne pas les injections et n’intervient pas dans la réduction des cylindres.
L’ordre causal demeure APP--TRIT--TPK, état enrichi, structure discrète
conjointe du continu, projection archimédienne et identification ultérieure.

Sont ainsi démontrées conjointement la construction du continu HMT, l'équipotence de sa droite avec l'ensemble des parties de \(\omega\), les deux injections cardinales, l'absence de cardinaux intermédiaires dans la clôture complète et l'égalité \(2^{\aleph_0}=\aleph_1\). Telle est la résolution affirmative de l'hypothèse du continu établie par le théorème \ref{thm:decision-continuo-hmt}.

\endgroup
